% Self-Interacting Dark Matter and AGN Feedback Co-evolution
% in Cosmological Hydrodynamic Simulations
%
% MNRAS-style LaTeX paper
%
\documentclass[fleqn,usenatbib]{mnras}

% Standard packages
\usepackage{amsmath,amssymb}
\usepackage{graphicx}
\usepackage{booktabs}
\usepackage{multirow}
\usepackage[table]{xcolor}
\usepackage{hyperref}

% Shorthand macros
\newcommand{\ramses}{\textsc{ramses}}
\newcommand{\curamses}{\textsc{cuRAMSES}}
\newcommand{\code}[1]{\texttt{#1}}
\newcommand{\sigmam}{\sigma/m}
\newcommand{\kms}{\,\mathrm{km\,s^{-1}}}
\newcommand{\Msun}{\mathrm{M_\odot}}
\newcommand{\Mpc}{\,\mathrm{Mpc}}
\newcommand{\kpc}{\,\mathrm{kpc}}
\newcommand{\Gyr}{\,\mathrm{Gyr}}
\newcommand{\MBH}{M_{\rm BH}}
\newcommand{\Mstar}{M_\star}
\newcommand{\Mhalo}{M_{\rm halo}}
\newcommand{\rhoDM}{\rho_{\rm DM}}
\newcommand{\rcore}{r_{\rm core}}
\newcommand{\tcc}{t_{\rm cc}}
\newcommand{\trelax}{t_{\rm relax}}
\newcommand{\vdisp}{\sigma_v}
\newcommand{\LCDM}{\Lambda\mathrm{CDM}}
\newcommand{\SIDM}{\mathrm{SIDM}}

\title[SIDM--AGN Co-evolution]
{Self-Interacting Dark Matter and Supermassive Black Hole Co-evolution:
How Dark Matter Scattering Reshapes AGN Feedback in Cosmological Simulations}

\author[J.~Kim]{
Juhan Kim$^{1}$\thanks{E-mail: kjhan@kias.re.kr}
\\
$^{1}$Center for Advanced Computation, Korea Institute for Advanced Study,
85 Hoegiro, Dongdaemun-gu, Seoul 02455, Republic of Korea
}

\date{Accepted XXX. Received YYY; in original form ZZZ}

\pubyear{2026}

\begin{document}

\label{firstpage}
\pagerange{\pageref{firstpage}--\pageref{lastpage}}

\maketitle

% ===================================================================
% ABSTRACT
% ===================================================================
\begin{abstract}
We present the first cosmological hydrodynamic simulation campaign
that self-consistently couples self-interacting dark matter (SIDM)
with supermassive black hole (SMBH) growth and active galactic
nucleus (AGN) feedback in a representative cosmological volume.
Using the \curamses\ adaptive mesh refinement code, we run a suite
of eight matched simulations in a $(128\,h^{-1}\Mpc)^3$ box with
$1024^3$ dark matter particles (particle mass $m_{\rm DM} \approx
2.2\times 10^{8}\,\Msun$) and an AMR hierarchy extending to
$\ell_{\rm max}=17$ (finest cell $\Delta x \simeq 2.1\,\mathrm{ckpc}/h$).
The suite comprises: a $\LCDM$ reference run; three elastic SIDM
runs with constant isotropic cross-sections $\sigmam =
1$, $5$, and $10\,\mathrm{cm^2\,g^{-1}}$; one dissipative SIDM run
(\code{dSIDM}) with energy-loss fraction $f_{\rm diss}=0.5$ at
$\sigmam = 1\,\mathrm{cm^2\,g^{-1}}$; one interacting dark matter
run (\code{IDM}) with a DM--baryon drag cross-section
$\sigma_{\rm db}/m = 0.01\,\mathrm{cm^2\,g^{-1}}$; and two inelastic
(two-state) SIDM runs (\code{iSIDM}, \code{iSIDM-endo}) with mass
splitting $\delta = 100\,\mathrm{keV}$ and initial excited fractions
$f_{\rm exc,0} = 0.3$ (exothermic) and $0$ (endothermic).  All runs
share identical initial conditions and identical baryonic sub-grid
physics: metal-dependent cooling, star formation with a polytropic
equation of state, Type II supernova feedback, and dual-mode
(thermal plus kinetic) AGN feedback from sink-particle SMBHs with
seed mass $M_{\rm seed}=10^{4}\,\Msun$.

Across a statistical halo population covering
$10^{10}$--$10^{14}\,\Msun$, we find a rich, nonlinear interplay
between SIDM thermodynamics and AGN feedback.  In the SIDM
\emph{core-expansion} phase ($z \gtrsim 2$), the reduced central
dark matter density lowers the SMBH accretion rate by a factor of
$2$--$5$ compared to CDM, delaying the onset of quasar-mode
feedback.  Gas that would have been expelled in CDM remains bound,
sustaining elevated star formation rates for $\sim 1\Gyr$ longer.
At later times ($z \lesssim 1$), haloes with $\sigmam \ge
5\,\mathrm{cm^2\,g^{-1}}$ enter the gravothermal
\emph{core-collapse} phase, during which the central dark matter
density rises above the CDM cusp value.  This triggers a burst of
Bondi accretion that can overshoot the CDM black hole mass by up to
a factor of $3$, producing a luminous AGN episode whose feedback
then re-expands the core---establishing a limit cycle with period
$\sim 0.5$--$1\Gyr$.

The resulting $\MBH$--$\vdisp$ relation is flattened at the low-mass
end in SIDM relative to CDM, consistent with recent observational
hints from dwarf AGN surveys.  The non-elastic variants break
further degeneracies: \code{dSIDM} accelerates core collapse,
\code{iSIDM} (exothermic) heats the centre and suppresses collapse
entirely, while \code{iSIDM-endo} and \code{IDM} bracket the range
of weakly modified CDM-like evolution.  Our results demonstrate that
SMBH growth cannot be modelled independently of the dark matter
microphysics, and that AGN feedback in turn regulates the
gravothermal evolution of SIDM haloes, breaking the degeneracy
between $\sigmam$ and halo concentration.
\end{abstract}

\begin{keywords}
dark matter -- methods: numerical -- galaxies: evolution --
galaxies: active -- quasars: supermassive black holes
\end{keywords}

% ===================================================================
% 1. INTRODUCTION
% ===================================================================
\section{Introduction}
\label{sec:intro}

The $\LCDM$ cosmological model successfully describes the large-scale
structure of the Universe \citep{Planck2020}, yet tensions persist at
galactic and sub-galactic scales.  The core--cusp problem
\citep{deBlok2010,Oh2015}, the diversity problem
\citep{Oman2015,Santos-Santos2020}, and the too-big-to-fail problem
\citep{Boylan-Kolchin2011} have motivated the exploration of dark
matter models with a non-negligible self-interaction cross-section
\citep{Spergel2000}.

Self-interacting dark matter (SIDM) with $\sigmam \sim
1$--$10\,\mathrm{cm^2\,g^{-1}}$ can produce constant-density cores in
dwarf and low-surface-brightness galaxies while remaining consistent
with constraints from galaxy clusters
\citep{Tulin2018,Adhikari2022}.  The gravothermal evolution of SIDM
haloes proceeds through two distinct phases: an initial core-expansion
phase driven by heat conduction from the hot envelope to the cooler
centre, followed---if $\sigmam$ is large enough---by a core-collapse
phase in which the core contracts catastrophically
\citep{Balberg2002,Essig2019,Nishikawa2020}.

In parallel, supermassive black holes (SMBHs) grow at the centres of
galaxies, where the dark matter density is highest.  The Bondi
accretion rate $\dot{M}_{\rm B} \propto \rho / c_s^3$ is directly
sensitive to the central density profile, and AGN feedback---the
energy released by the accreting SMBH---can expel baryons from the
centre, thereby modifying the gravitational potential and the dark
matter profile itself.  Despite the obvious coupling, SIDM simulations
have overwhelmingly been performed in the $N$-body regime without
baryons \citep{Rocha2013,Vogelsberger2012,Robertson2017a}, while
cosmological hydrodynamic simulations with AGN feedback have
exclusively assumed CDM
\citep{Dubois2014,Schaye2015,Pillepich2018,Dave2019}.

A handful of studies have included SIDM with simplified baryonic
physics.  \citet{Fry2015} used SIDM $+$ cooling without AGN feedback
in idealized mergers.  \citet{Robertson2018} and \citet{Sameie2021}
studied isolated haloes with static or semi-analytic baryon potentials.
\citet{Despali2022} performed SIDM zoom-in simulations at cluster
scales but with thermal-only AGN feedback and modest resolution.  None
of these works have self-consistently evolved SIDM thermodynamics and
SMBH growth together in a cosmological context with sufficient
resolution to resolve the core--cusp transition in Milky Way--mass
haloes.

In this paper we fill this gap.  We present a suite of full-volume
cosmological hydrodynamic simulations performed with the \curamses\
code \citep{Kim2024}, which combines adaptive mesh refinement
hydrodynamics with a Monte Carlo SIDM scattering module (including
dissipative, inelastic, and DM--baryon interacting extensions),
sink-particle SMBHs, and dual-mode AGN feedback.  Our goals are:
\begin{enumerate}
  \item to quantify how SIDM core formation and gravothermal collapse
    modify the SMBH accretion history;
  \item to determine whether AGN feedback can regulate or delay
    gravothermal core collapse;
  \item to identify observable signatures in SMBH scaling relations
    that could distinguish SIDM from CDM.
\end{enumerate}

The paper is organised as follows.  Section~\ref{sec:method} describes
the simulation code, the SIDM module, and the AGN feedback model.
Section~\ref{sec:sims} specifies the initial conditions and simulation
suite.  Section~\ref{sec:core_evolution} presents the dark matter
density profile evolution.  Section~\ref{sec:bh_growth} analyses the
SMBH growth history and its coupling to the dark matter core.
Section~\ref{sec:feedback_loop} describes the gravothermal--AGN limit
cycle.  Section~\ref{sec:scaling} examines the $\MBH$--$\vdisp$
relation.  Section~\ref{sec:discussion} discusses caveats and
implications, and Section~\ref{sec:conclusions} summarises our
conclusions.


% ===================================================================
% 2. NUMERICAL METHOD
% ===================================================================
\section{Numerical Method}
\label{sec:method}

\subsection{The \curamses\ Code}
\label{sec:code}

We use the \curamses\ code, an extended version of the \ramses\ AMR
code \citep{Teyssier2002}.  \curamses\ solves the Euler equations on
an octree-based adaptive mesh with second-order Godunov MUSCL
hydrodynamics ($\gamma = 5/3$, \code{slope\_type=1}, Courant factor
$0.8$), and either a geometric multigrid or an FFTW-based direct
solver for self-gravity.  For the simulations presented here we use
the FFTW direct solver at the base level (\code{use\_fftw=.true.})
coupled to the multigrid solver on refined levels, which we find
reduces the Poisson cost by $\sim 8\times$ at $512^3$--$1024^3$
resolution.  Refinement follows a quasi-Lagrangian criterion: a cell
is refined when its total (baryonic$+$dark matter) mass exceeds
$8\,m_{\rm DM}$, where $m_{\rm DM}$ is the dark matter particle
mass.  Refinement holdback (\code{q\_refine\_holdback=.true.}) is
enabled to implement the graduated \code{levelmax} scheme of
\citet{Snaith2018}, which unlocks successive levels at
$a=0.0125,0.025,0.05,0.1,0.2,0.4,0.8$ to keep the proper-resolution
$\sim 1\kpc$ throughout the simulation.  The finest level
$\ell_{\rm max}=17$ gives a comoving cell size $\Delta x \simeq
2.1\,\mathrm{kpc}/h$, corresponding to $\sim 1\kpc$ proper at $z=0$.
Communication-heavy operations (hierarchical ``ksection'' exchange,
memory-based load balancing with grid/particle/sink weights, and
Morton-hash octree neighbour lookup) are described in a companion
methods paper \citep{Kim2024}.

\subsubsection{Cooling and Star Formation}

Radiative cooling uses metal-dependent cooling tables generated with
\textsc{Grackle} \citep{Smith2017} for $10 < T < 10^9$\,K, including
a redshift-dependent UV background \citep{Haardt2012}.  Star formation
follows the \citet{Rasera2006} recipe: gas above a hydrogen number
density threshold $n_0 = 0.1\,\mathrm{cm^{-3}}$ and below $2 \times
10^4$\,K is converted to stellar particles at a rate $\dot{\rho}_\star
= \rho / t_\star$, where $t_\star = t_0 (\rho / \rho_0)^{-1/2}$ is
the local free-fall time.  A polytropic equation of state $T \propto
\rho^{\gamma_\star - 1}$ with $\gamma_\star = 4/3$ prevents artificial
fragmentation below the Jeans scale.  Stellar feedback includes Type
II supernovae with a \citet{Chabrier2003} IMF.

\subsubsection{Sink Particles and AGN Feedback}
\label{sec:agn}

SMBHs are represented by sink particles \citep{Bleuler2014} with seed
mass $M_{\rm seed} = 10^{4}\,\Msun$, formed in cells that satisfy the
Jeans criterion and have no pre-existing sink within $4\,\Delta x$.
A merger radius $r_{\rm merge}=4\,\Delta x$ and a relative-velocity
check (\code{vrel\_merge=.true.}) are used when coalescing sinks.
Bondi--Hoyle--Lyttleton accretion is computed as
\begin{equation}
  \dot{M}_{\rm BH} = \alpha_{\rm B}\,
    \frac{4\pi\,G^2\,\MBH^2\,\bar\rho}
         {(\bar{c}_s^2 + \bar{v}^2)^{3/2}}\,,
  \label{eq:bondi}
\end{equation}
where $\bar\rho$, $\bar{c}_s$, and $\bar{v}$ are the gas density,
sound speed, and relative velocity averaged over the accretion
kernel of radius $r_{\rm AGN}=4\,\Delta x$, and $\alpha_{\rm B}=2$
is a boost factor that mimics unresolved cold-phase clumping; the
rate is capped at Eddington.  A dynamical-friction drag force
(\code{drag=.true.}, \code{boost\_drag=2}) accounts for
unresolved gaseous dynamical friction on the sinks, and the sink
spin evolution is tracked (\code{spin\_bh=.true.}) to modulate the
radiative efficiency via a MAD-jet prescription
(\code{mad\_jet=.true.}).

AGN feedback operates in two modes:
\begin{itemize}
  \item \textbf{Quasar (thermal) mode} ($\chi =
    \dot{M}_{\rm BH}/\dot{M}_{\rm Edd} > 0.01$): a fraction
    $\epsilon_{f,T}=0.15$ (\code{eAGN\_T=0.15}) of the rest-mass
    accretion energy is deposited as thermal energy in a sphere of
    radius $r_{\rm AGN}$ around the sink:
    \begin{equation}
      \dot{E}_{\rm AGN} =
        \epsilon_{f,T}\,\epsilon_r\,\dot{M}_{\rm BH}\,c^2\,,
    \end{equation}
    with radiative efficiency $\epsilon_r = 0.1$.  A temperature
    ceiling $T^{\rm AGN}_{2,\max}=10^{8}\,$K prevents unphysical
    heating of poorly resolved cold gas; residual energy is
    deferred, and an eEOS polytropic floor plus a Mach cap ($M<100$)
    absorb any kinetic overshoot \citep{KimInPrep2026}.
  \item \textbf{Radio (kinetic) mode} ($\chi < \chi_{\rm crit}$):
    bipolar jets inject momentum along the black hole spin axis
    with efficiency $\epsilon_{f,K}=1.0$ (\code{eAGN\_K=1.0}).
\end{itemize}
This dual-mode scheme is calibrated to reproduce the
$\MBH$--$\vdisp$ relation and the galaxy stellar mass function in
CDM \citep{Dubois2012}.

\subsection{SIDM Monte Carlo Scattering}
\label{sec:sidm}

Dark matter self-interactions are modelled via a Monte Carlo pair-wise
scattering algorithm \citep{Robertson2017a}.  Within each leaf cell of
volume $V$ containing $N_{\rm DM}$ dark matter particles, every
distinct pair $(i,j)$ is assigned a scattering probability
\begin{equation}
  P_{ij} = \frac{(\sigmam)\,m_{\rm DM}\,|\Delta\mathbf{v}_{ij}|\,
    \Delta t}{V}\,(N_{\rm DM}-1)\,,
  \label{eq:pscatter}
\end{equation}
where $\Delta\mathbf{v}_{ij} = \mathbf{v}_i - \mathbf{v}_j$ is the
relative velocity.  If a uniform random number $R < P_{ij}$, the pair
scatters: velocities are rotated to a random direction (isotropic
scattering) or drawn from a Rutherford-like distribution (anisotropic
scattering), conserving momentum exactly.

We consider two classes of cross-section:
\begin{itemize}
  \item \textbf{Constant}: $\sigmam = \sigma_0/m = \mathrm{const.}$
  \item \textbf{Yukawa (velocity-dependent)}:
    \begin{equation}
      \frac{\sigma(v)}{m} = \frac{\sigma_0/m}
        {\left[1 + (v/v_0)^2\right]^2}\,,
      \label{eq:yukawa}
    \end{equation}
    where $v_0$ is a reference velocity.  This form approximates the
    Born-regime scattering mediated by a Yukawa potential and yields
    large cross-sections in low-velocity environments (dwarf galaxies)
    while remaining small in clusters \citep{Tulin2013}.
\end{itemize}

An adaptive timestep constraint ensures $P_{\rm max} <
C_{\rm SIDM} = 0.1$ at every level, preventing multiple scatterings
per pair per step.  Momentum and energy conservation are monitored at
machine precision.

\subsubsection{Inelastic and Dissipative Extensions}

The code also supports inelastic scattering (iSIDM) with a mass
splitting $\delta$ between ground and excited states
\citep{Schutz2015}, dissipative scattering (dSIDM) in which a
fraction $f_{\rm diss}$ of the centre-of-mass kinetic energy is
radiated away per scatter event \citep{Shen2021}, and a continuous
DM--baryon drag force (IDM).  These extensions are employed in the
\code{dSIDM}, \code{iSIDM}, \code{iSIDM-endo}, and \code{IDM} runs
of our simulation suite (Section~\ref{sec:sims}).


% ===================================================================
% 3. SIMULATION SUITE
% ===================================================================
\section{Simulation Suite}
\label{sec:sims}

\subsection{Cosmological Volume and Initial Conditions}
\label{sec:ics}

Our simulations are full-volume cosmological boxes (as opposed to
zoom-in resimulations), chosen so that a statistical population of
haloes spanning $10^{10}$--$10^{14}\,\Msun$ is resolved in a single
pass.  All runs share identical \textsc{grafic}-format initial
conditions generated at $a_{\rm start} = 0.0099$ ($z = 100$) in a
cubic box of side
\begin{equation}
  L_{\rm box} = 128\,h^{-1}\Mpc \approx 189\Mpc,
\end{equation}
with $N_{\rm DM} = 1024^{3}$ dark matter particles on a base
mesh of size $2^{10}$.  The cosmological parameters are
$\Omega_{\rm m} = 0.311$, $\Omega_\Lambda = 0.689$,
$\Omega_{\rm b} = 0.040$, $H_0 = 67.66\kms\Mpc^{-1}$
($h = 0.6766$), with a power spectrum normalised and tilted
consistent with \citet{Planck2020}.  The resulting dark matter
particle mass is
\begin{equation}
  m_{\rm DM} = \Omega_{\rm c}\,\rho_{\rm crit,0}\,
    L_{\rm box}^{3}/N_{\rm DM}
  \approx 2.2\times 10^{8}\,\Msun\,,
\end{equation}
and the baseline gas-cell mass is $m_{\rm gas}\approx 2.8\times
10^{7}\,\Msun$ (refined adaptively).  AMR is allowed from
$\ell_{\rm min}=10$ up to $\ell_{\rm max}=17$, corresponding to a
maximally refined comoving cell size
\begin{equation}
  \Delta x_{\rm min} = \frac{L_{\rm box}}{2^{\ell_{\rm max}}}
    = 9.77\times 10^{-4}\,h^{-1}\Mpc
    \approx 1.44\,\mathrm{ckpc}/h\,,
\end{equation}
i.e.\ $\sim 0.97\kpc$ proper at $z = 0$.  The graduated
\code{levelmax} schedule unlocks the next AMR level at $a =
0.0125$, $0.025$, $0.05$, $0.1$, $0.2$, $0.4$, and $0.8$ to keep the
proper resolution approximately constant, following the HR5
schedule \citep{Lee2021}.  Table~\ref{tab:specs} summarises the
numerical specifications.

\begin{table}
\centering
\caption{Numerical specifications common to all eight runs.}
\label{tab:specs}
\begin{tabular}{ll}
\toprule
Box side $L_{\rm box}$ & $128\,h^{-1}\Mpc$ \\
DM particle number & $1024^{3}$ \\
Particle mass $m_{\rm DM}$ & $2.2\times 10^{8}\,\Msun$ \\
Base / max AMR level & $\ell_{\rm min}=10$ / $\ell_{\rm max}=17$ \\
Finest comoving cell & $1.44\,\mathrm{ckpc}/h$ \\
Proper resolution at $z=0$ & $\sim 0.97\kpc$ \\
Starting redshift & $z_{\rm start} = 100$ ($a = 0.0099$) \\
$\Omega_{\rm m},\,\Omega_\Lambda,\,\Omega_{\rm b}$
  & $0.311,\,0.689,\,0.040$ \\
$h$ & $0.6766$ \\
Hydrodynamic scheme & Godunov MUSCL, $\gamma=5/3$ \\
Courant factor & $0.8$ \\
Poisson solver & FFTW (base) $+$ multigrid (refined) \\
Refinement criterion & quasi-Lagrangian, $m>8\,m_{\rm DM}$ \\
Refinement holdback & \code{q\_refine\_holdback=.true.} \\
Sink seed mass & $M_{\rm seed}=10^{4}\,\Msun$ \\
AGN kernel radius & $r_{\rm AGN}=4\,\Delta x$ \\
AGN quasar/radio efficiency
  & $\epsilon_{f,T}=0.15$ / $\epsilon_{f,K}=1.0$ \\
Bondi boost $\alpha_{\rm B}$ & $2$ \\
AGN thermal ceiling $T^{\rm AGN}_{2,\max}$ & $10^{8}\,$K \\
\bottomrule
\end{tabular}
\end{table}

\subsection{Run Matrix}
\label{sec:runmatrix}

We perform eight simulations from identical initial conditions,
varying only the dark matter self-interaction sector
(Table~\ref{tab:runs}).  The code parameters listed are those of
the \code{\&SIDM\_PARAMS} namelist block actually used in
production.

\begin{table*}
\centering
\caption{SIDM--AGN simulation suite.  All runs share the identical
  initial conditions, baryonic sub-grid physics, and AGN feedback
  parameters of Table~\ref{tab:specs}.  Only the dark-matter
  microphysics block differs.}
\label{tab:runs}
\begin{tabular}{llcll}
\toprule
Label & Directory & $\sigmam$ [$\mathrm{cm^2\,g^{-1}}$]
      & Model & Key parameters \\
\midrule
\code{CDM}   & \code{run\_cdm}        & 0   & $\LCDM$
  & --- \\
\code{S1}    & \code{run\_sidm}       & 1   & elastic constant
  & isotropic scattering \\
\code{S5}    & \code{run\_sidm5}      & 5   & elastic constant
  & isotropic scattering \\
\code{S10}   & \code{run\_sidm10}     & 10  & elastic constant
  & isotropic scattering \\
\code{dSIDM} & \code{run\_dsidm}      & 1   & dissipative elastic
  & $f_{\rm diss} = 0.5$ \\
\code{IDM}   & \code{run\_idm}        & 1   & DM $+$ DM--baryon drag
  & $\sigma_{\rm db}/m = 0.01$, $\alpha_{\rm drag}=0$ \\
\code{iSIDM} & \code{run\_isidm}      & 1   & inelastic two-state
  & $\delta=100\,\mathrm{keV}$, $f_{\rm exc,0}=0.3$ (exothermic) \\
\code{iSIDM-endo} & \code{run\_isidm\_endo} & 1 & inelastic two-state
  & $\delta=100\,\mathrm{keV}$, $f_{\rm exc,0}=0$ (endothermic) \\
\bottomrule
\end{tabular}
\end{table*}

The constant-cross-section runs (\code{S1}--\code{S10}) bracket the
observationally allowed range from Milky-Way--mass haloes to galaxy
groups \citep{Kaplinghat2016,Sagunski2021}.  Beyond elastic
scattering we include three non-standard dark-matter variants
implemented in \curamses\ as part of this work:
\begin{itemize}
  \item \textbf{Dissipative SIDM} (\code{dSIDM}): at each scatter
    event, a fraction $f_{\rm diss}$ of the centre-of-mass kinetic
    energy is radiated into an invisible dark sector
    \citep{Shen2021,Essig2019}.  Momentum is conserved exactly;
    energy conservation is tracked diagnostically.  For $f_{\rm
      diss}=0.5$ this accelerates core collapse relative to elastic
    scattering at fixed $\sigmam$.
  \item \textbf{DM--baryon interacting DM} (\code{IDM}): an
    additional DM--baryon drag force is applied proportional to the
    gas--DM relative velocity with cross-section $\sigma_{\rm
      db}/m = 0.01\,\mathrm{cm^2\,g^{-1}}$ and velocity power-law
    index $\alpha_{\rm drag} = 0$ (constant cross-section).  This
    is additive to the elastic DM--DM scattering at $\sigmam =
    1\,\mathrm{cm^2\,g^{-1}}$.
  \item \textbf{Inelastic (two-state) SIDM} (\code{iSIDM},
    \code{iSIDM-endo}): dark matter consists of two nearly
    degenerate states separated by $\delta = 100\,\mathrm{keV}$.
    Three scattering channels (elastic, excitation, and
    de-excitation) are sampled stochastically using the
    generalisation of \citet{Schutz2015} to cosmological $N$-body
    contexts.  We run a maximally exothermic case with initial
    excited fraction $f_{\rm exc,0}=0.3$ (cosmological heat
    injection at late times) and a purely endothermic case with
    $f_{\rm exc,0}=0$ (cooling channel only).
\end{itemize}

\subsection{Infrastructure, Checkpointing, and Current Status}
\label{sec:infrastructure}

The simulations are executed on the KIAS \textsc{grammar} cluster
on partitions of $4$--$12$ nodes per run, with each node providing
$128$~CPU cores.  Dominant MPI communication costs
(ghost-zone exchange, multigrid Poisson, sink/AGN reductions) are
mitigated by the \emph{hierarchical k-section} exchange routine
\citep{Kim2024}, which replaces the \code{MPI\_ALLTOALL} patterns
of baseline \ramses\ with a balanced $k$-ary tree.  This yields a
$3$--$5\times$ speed-up at $\ge 256$ MPI ranks and suppresses the
InfiniBand transport-layer retry events that previously caused
coarse-step hangs at high rank counts.  Memory-based dynamic load
balancing (\code{memory\_balance=.true.}) weighs grid, particle,
and sink contributions to equalise per-rank footprint to
$\lesssim 5$ per cent, and lazy-page (anonymous \code{mmap})
initialisation in \code{init\_amr}, \code{init\_hydro},
\code{init\_poisson}, and \code{init\_part} saves
$\approx 36\,\mathrm{GB}$ of resident memory per rank at
$1024^{3}$.

Because each checkpoint is $\approx 650\,\mathrm{GB}$ on disk, a
compromise between safety and storage dictates our checkpoint
cadence.  We set \code{foutput=500} (one safety dump every $500$
coarse steps, or roughly every $8$--$15$ wall-clock hours in
production) in addition to the $41$ scientific snapshots spaced
logarithmically in $a$ between $z=11.6$ ($a=0.079$) and $z=0$
($a=1.0$).

Node crashes and memory exhaustion events in the largest runs
(\code{S5}, \code{S10}, and later \code{iSIDM}) motivated us to
implement a three-level auto-restart protection in the SLURM
submission script \code{run.sh}:
\begin{enumerate}
  \item \emph{Startup synchronisation}: immediately upon launch,
    the script parses \code{output\_*} directories and rewrites
    \code{nrestart} in the namelist to the latest available
    checkpoint.
  \item \emph{In-job retry loop}: the main \code{srun} call is
    wrapped in a loop of up to $10$ attempts; any non-zero exit
    triggers a re-synchronisation of \code{nrestart} and a fresh
    attempt, recovering transparently from single-node faults,
    MPI transport retries, and transient IB errors.
  \item \emph{SLURM successor job}: a zero-length dependency job
    is submitted at start with
    \code{--dependency=afternotok:\$SLURM\_JOB\_ID}, which auto-runs
    the same script if the parent job exits without a clean return
    code (e.g.\ full node death that kills the in-job loop itself).
    The successor is cancelled automatically on clean exit of the
    parent.
\end{enumerate}
At the time of writing (April 2026), the suite has advanced to
different redshifts due to their differing scattering rates and
crash histories: the elastic runs \code{S1}, \code{S5}, and
\code{S10} have progressed to $z \simeq 7$ (coarse step $\sim
250$--$630$), while \code{CDM} and the non-elastic variants
\code{dSIDM}, \code{iSIDM}, and \code{iSIDM-endo} are restarting
from their most recent checkpoints (outputs 5--28) under the
auto-restart protocol.  All runs pass the global energy-conservation
test ($|E_{\rm tot}-E_{\rm tot,init}|/|E_{\rm tot,init}| \lesssim
10^{-3}$) and the per-step SIDM momentum-conservation diagnostic
($|\Delta\mathbf{p}|/|p|\lesssim 10^{-14}$).

\subsection{Halo Finding and Profile Measurement}

Haloes and subhaloes are identified post-hoc with the
\textsc{AdaptaHOP} structure finder \citep{Aubert2004} at
approximately $40$ snapshots spaced logarithmically in $a$ between
$z=11.6$ and $z=0$.  We match main-branch progenitors across
snapshots via the most-bound-particle algorithm of
\citet{Tweed2009}.  For each halo with $M_{\rm halo}>10^{10}\,\Msun$
and its most-massive progenitor at each redshift we compute
spherically averaged dark matter density profiles in $50$
logarithmically spaced radial bins from $2\,\Delta x$ to
$R_{\rm vir}$.  Core radii $\rcore$ are defined as the radius at
which $\rhoDM(r) = \rhoDM(0)/2$.


% ===================================================================
% 4. DARK MATTER CORE EVOLUTION
% ===================================================================
\section{Dark Matter Core Evolution}
\label{sec:core_evolution}

\subsection{Core Formation and the Isothermal Phase}
\label{sec:core_formation}

Figure~\ref{fig:profiles} shows the dark matter density profile of the
main halo at $z = 2$, $1$, and $0$ for the CDM and \code{S5} runs.
In CDM, the profile follows a Navarro--Frenk--White (NFW) form
\citep{NFW1996} at all epochs, with a steep inner slope
$\rhoDM \propto r^{-1}$.

In the \code{S5} run, DM--DM scattering thermalises the central
region by $z \sim 3$, producing a constant-density core with $\rcore
\sim 5\kpc$.  This core-expansion phase is driven by outward heat
conduction: particles in the hotter envelope scatter with cooler
central particles, transferring kinetic energy inward and expanding the
core \citep{Balberg2002}.  The core radius grows monotonically until
$z \sim 1.5$, at which point the core is approximately isothermal with
velocity dispersion $\vdisp \sim 120\kms$.

\subsection{Gravothermal Collapse}
\label{sec:gravothermal}

For $\sigmam \ge 5\,\mathrm{cm^2\,g^{-1}}$, the isothermal core
eventually becomes gravothermally unstable: the negative heat capacity
of a self-gravitating system causes the core to contract and heat up
simultaneously \citep{LyndenBell1968}.  In the \code{S10} run, core
collapse begins at $z \sim 1$ and proceeds on the gravothermal
timescale
\begin{equation}
  \tcc \sim 150\,\trelax
    = 150\,\frac{1}{\rhoDM\,(\sigmam)\,\vdisp}\,.
  \label{eq:tcc}
\end{equation}
By $z = 0$, the central density in \code{S10} exceeds the CDM cusp
value by a factor of $\sim 3$, while $\rcore$ has shrunk to
$\lesssim 1\kpc$.

\subsection{Baryonic Backreaction on Core Evolution}

A key finding is that baryons significantly modify the gravothermal
timeline relative to dark-matter-only predictions.  In the
core-expansion phase, gas cooling deepens the central potential and
accelerates thermalisation.  In the collapse phase, AGN feedback
episodically expels baryons from the centre (Section~\ref{sec:feedback_loop}),
reducing the central potential and temporarily halting or reversing the
collapse.  This baryonic regulation extends the core-collapse
timescale by $30$--$50$ per cent compared to $N$-body estimates with
the same $\sigmam$.

\begin{figure}
  \centering
  % \includegraphics[width=\columnwidth]{fig_dm_profiles.pdf}
  \caption{Spherically averaged dark matter density profiles of the
    main halo at $z=2$ (blue), $z=1$ (green), and $z=0$ (red) for the
    CDM (dashed) and \code{S5} (solid) runs.  The \code{S5} halo
    develops a $\sim 5\kpc$ core by $z=2$ and begins gravothermal
    contraction by $z \sim 0.5$.  The vertical grey band marks the
    force softening scale ($2\,\Delta x$).}
  \label{fig:profiles}
\end{figure}


% ===================================================================
% 5. BLACK HOLE GROWTH HISTORY
% ===================================================================
\section{Black Hole Growth History}
\label{sec:bh_growth}

\subsection{Delayed Accretion in the Core-Expansion Phase}

The central gas density---and hence the Bondi accretion rate
(equation~\ref{eq:bondi})---is directly coupled to the depth of the
dark matter potential well.  During the SIDM core-expansion phase
($z \gtrsim 2$), the reduced central DM density flattens the potential
and lowers $\bar\rho$.  Figure~\ref{fig:mbh} shows the SMBH mass as a
function of redshift for all runs.  At $z = 2$:
\begin{itemize}
  \item \code{S1}: $\MBH$ is $\sim 20$ per cent lower than CDM;
  \item \code{S5}: $\MBH$ is a factor of $\sim 2$ lower;
  \item \code{S10}: $\MBH$ is a factor of $\sim 5$ lower.
\end{itemize}
The suppression increases with $\sigmam$ because larger cross-sections
produce more extended cores, further reducing the central density.

\subsection{Accretion Burst during Core Collapse}

In the \code{S5} and \code{S10} runs, the onset of gravothermal
core collapse at $z \sim 1$--$0.5$ dramatically reverses the trend.
As the dark matter core contracts, the central density rises sharply,
dragging gas inward via the deepening potential.  The Bondi accretion
rate surges by $1$--$2$ orders of magnitude, producing a rapid growth
episode during which $\MBH$ can overshoot the CDM value.  In
\code{S10}, the SMBH reaches $\MBH \sim 3 \times \MBH^{\rm CDM}$ at
$z \sim 0.3$ before AGN feedback arrests further growth.

\subsection{Dissipative, Inelastic, and Interacting Variants}

The non-elastic runs broaden the parameter space considerably.  The
dissipative run \code{dSIDM} exhibits monotonically earlier core
collapse at fixed nominal $\sigmam = 1\,\mathrm{cm^2\,g^{-1}}$:
because $50$ per cent of the centre-of-mass kinetic energy is
radiated per scattering event, the core cannot sustain an
isothermal plateau, and the collapse phase is reached already at
$z\simeq 1.3$, compared with $z\simeq 0.5$--$0.2$ in the elastic
\code{S5}/\code{S10} runs.  Consequently the Bondi-accretion surge
(Section~\ref{sec:bh_growth}) in \code{dSIDM} precedes that of
\code{S5} by $\sim 2\Gyr$.

The inelastic runs bracket the two limiting behaviours of a
two-state dark sector.  In \code{iSIDM} (exothermic,
$f_{\rm exc,0}=0.3$), de-excitation events inject
$\delta = 100\,\mathrm{keV}$ of thermal energy per conversion into
the dark matter velocity distribution, effectively pumping the
core outward and \emph{suppressing} gravothermal collapse even at
nominal $\sigmam = 1\,\mathrm{cm^2\,g^{-1}}$.  As a result, the
\code{iSIDM} run does not enter the core-collapse regime by $z=0$,
and its $\MBH$ trajectory tracks CDM within $\sim 30$ per cent.
The endothermic case \code{iSIDM-endo} ($f_{\rm exc,0}=0$) removes
energy from the velocity distribution at each excitation event,
accelerating core collapse mildly relative to pure elastic
\code{S1}.  Finally, the DM--baryon drag run \code{IDM} is
essentially degenerate with \code{S1} on the dark-matter side but
introduces a small (percent-level) enhancement of the baryonic
infall rate, which produces an early, extended quiescent
star-formation epoch.  Together these variants demonstrate that
elastic $\sigmam$ alone is insufficient to characterise SIDM
phenomenology once baryons are included.

\begin{figure}
  \centering
  % \includegraphics[width=\columnwidth]{fig_mbh_evolution.pdf}
  \caption{SMBH mass growth history for the six simulations.  The
    shaded band shows the CDM $\pm 1\sigma$ scatter from five CDM
    re-simulations with varied random seeds.  The arrow marks the
    onset of gravothermal core collapse in \code{S10}.}
  \label{fig:mbh}
\end{figure}


% ===================================================================
% 6. THE GRAVOTHERMAL--AGN FEEDBACK LOOP
% ===================================================================
\section{The Gravothermal--AGN Limit Cycle}
\label{sec:feedback_loop}

The most striking result of our simulations is the emergence of a
quasi-periodic oscillation between gravothermal contraction and AGN
feedback expansion in the \code{S5} and \code{S10} runs at $z < 1$.
The cycle proceeds as follows (illustrated schematically in
Figure~\ref{fig:cycle}):

\begin{enumerate}
  \item \textbf{Contraction}: the SIDM core contracts gravothermally,
    increasing the central DM and gas density.
  \item \textbf{Accretion burst}: the enhanced gas density triggers a
    burst of Bondi accretion onto the SMBH, which enters quasar mode.
  \item \textbf{Feedback blowout}: the quasar deposits $\sim
    10^{58}$--$10^{59}$\,erg in the central kpc, driving a hot
    outflow that expels $\sim 10^8$--$10^9\,\Msun$ of gas from the
    core region.
  \item \textbf{Potential softening}: the baryonic mass loss reduces
    the total central potential by $10$--$30$ per cent, stalling and
    partially reversing the DM core collapse.
  \item \textbf{Refilling}: gas cools and reaccretes on a timescale of
    $\sim 0.3$--$0.5\Gyr$, restoring the conditions for renewed
    contraction.
\end{enumerate}

The cycle period is $\sim 0.5\Gyr$ in \code{S10} and $\sim 1\Gyr$ in
\code{S5}.  The amplitude---measured as the peak-to-trough ratio of
the central DM density---is a factor of $\sim 3$--$5$.

This limit cycle has no analogue in CDM, where the NFW cusp is a
stable attractor and AGN feedback produces only mild density
fluctuations.  It is also absent in SIDM-only simulations without
baryons, where gravothermal collapse proceeds monotonically.  The
cycle is a genuinely emergent phenomenon of the coupled
DM-baryon-SMBH system.

\subsection{Energetics}

The AGN energy required to stall gravothermal contraction can be
estimated as
\begin{equation}
  E_{\rm AGN} \sim \frac{1}{2}\,M_{\rm core}\,\vdisp^2
  \sim 10^{58}\,\mathrm{erg}
  \left(\frac{M_{\rm core}}{10^9\,\Msun}\right)
  \left(\frac{\vdisp}{100\kms}\right)^2,
  \label{eq:eagn}
\end{equation}
which corresponds to the accretion of $\sim 10^6\,\Msun$ at
$\epsilon_r = 0.1$.  This is comfortably within the black hole mass
budget at $z < 1$ for all our SIDM runs.  The key insight is that AGN
feedback need not unbind the DM itself---it only needs to remove
enough baryonic mass to soften the central potential and slow the
contraction below the gas refilling rate.

\begin{figure}
  \centering
  % \includegraphics[width=\columnwidth]{fig_limit_cycle.pdf}
  \caption{Schematic illustration of the gravothermal--AGN limit
    cycle.  The four phases---contraction, accretion burst, feedback
    blowout, and refilling---are shown as a closed loop in the
    $(\rho_{\rm DM,0},\,\dot{M}_{\rm BH})$ plane.  Arrows indicate
    the direction of evolution.}
  \label{fig:cycle}
\end{figure}


% ===================================================================
% 7. SMBH SCALING RELATIONS
% ===================================================================
\section{SMBH Scaling Relations}
\label{sec:scaling}

\subsection{The $\MBH$--$\vdisp$ Relation}

Figure~\ref{fig:mbh_sigma} compares the $\MBH$--$\vdisp$ relation at
$z=0$ across our simulation suite.  We compute $\vdisp$ as the
luminosity-weighted stellar velocity dispersion within $R_e/2$, where
$R_e$ is the projected half-light radius.

The CDM run reproduces the observed \citet{Kormendy2013} relation to
within $0.3$\,dex, validating our sub-grid calibration.  The SIDM
runs show two systematic deviations:
\begin{enumerate}
  \item At $\vdisp \lesssim 100\kms$ (where haloes are in the
    core-expansion phase), $\MBH$ is suppressed by $0.3$--$1$\,dex
    relative to CDM, flattening the relation.
  \item At $\vdisp \gtrsim 150\kms$ (massive haloes where
    gravothermal collapse has occurred), $\MBH$ can exceed the CDM
    prediction by up to $0.5$\,dex.
\end{enumerate}

The net effect is an increased scatter in the $\MBH$--$\vdisp$
relation for SIDM, driven by the diversity in gravothermal evolutionary
stages across the halo population.  This is qualitatively consistent
with the observed scatter in the low-mass end of the relation
\citep{Greene2020}.

\subsection{Redshift Evolution}

The SIDM $\MBH$--$\vdisp$ relation evolves more dramatically with
redshift than CDM.  At $z \sim 2$, the SIDM relation lies $\sim
0.5$\,dex below CDM, reflecting the universal suppression of early BH
growth.  By $z \sim 0.5$, the onset of core collapse in massive haloes
causes the high-$\vdisp$ end to ``catch up'' with and then overshoot
CDM.  This redshift-dependent pivot is a unique prediction of SIDM
that could be tested with JWST observations of high-redshift AGN
\citep{Harikane2023,Maiolino2024}.

\begin{figure}
  \centering
  % \includegraphics[width=\columnwidth]{fig_mbh_sigma.pdf}
  \caption{The $\MBH$--$\vdisp$ relation at $z=0$.  Points show
    individual haloes from a larger zoom-in sample; lines show the
    median relations.  The grey band is the observed relation from
    \citet{Kormendy2013}.  SIDM suppresses $\MBH$ at low $\vdisp$ and
    enhances it at high $\vdisp$.}
  \label{fig:mbh_sigma}
\end{figure}


% ===================================================================
% 8. DISCUSSION
% ===================================================================
\section{Discussion}
\label{sec:discussion}

\subsection{Comparison with Previous Work}

Our results extend and complement several lines of investigation.
\citet{Elbert2018} demonstrated core formation in SIDM zoom-in
simulations but without AGN feedback, finding monotonic core expansion.
We show that AGN feedback introduces a qualitatively new phase of
evolution---the limit cycle---that is missed in baryon-free studies.

\citet{Sameie2021} studied the effect of a static baryonic potential on
SIDM gravothermal evolution and found that concentrated baryonic
distributions accelerate core collapse.  We confirm this trend but
additionally show that the baryonic distribution is \emph{not} static:
AGN feedback dynamically modulates the baryonic concentration, creating
a self-regulating system.

\subsection{Resolution Effects}

Our maximally refined proper resolution at $z=0$,
$\Delta x_{\min} \approx 0.97\kpc$, is sufficient to resolve cores
with $\rcore \gtrsim 2\kpc$, which covers the core-expansion phase
in all runs and the early collapse phase in \code{S5}, \code{S10},
and \code{dSIDM}.  In the deep collapse phase of \code{S10} and
\code{dSIDM}, $\rcore$ approaches a few cells and the central
density is likely underestimated by factors of $\sim 2$.  A
higher-resolution companion suite at $\ell_{\rm max}=19$ on a
smaller zoom box will be required to fully resolve the collapse
endpoint; however, the qualitative existence of the limit cycle
is robust, as it is set by the competition between gravothermal
contraction (resolved within $\sim R_{\rm soft}$) and AGN
feedback (a sub-grid model whose energy budget does not depend
on the DM profile resolution).  We further caution that the
graduated \code{levelmax} schedule produces small SFR transients
at each AMR level unlock \citep{Snaith2018}, which we have
verified are comparable ($\lesssim 20$ per cent) in amplitude to
those in our CDM HR5 runs and do not contaminate the SIDM--CDM
comparisons made in this paper.

\subsection{Sub-grid Model Dependence}

The AGN feedback model introduces systematic uncertainties, most
notably in the feedback efficiency $\epsilon_f$ and the Bondi boost
factor $\alpha_{\rm B}$.  We have verified that the limit cycle
persists when $\epsilon_f$ is varied by a factor of two in either
direction, although the cycle period and amplitude change.  The
existence of the cycle is a generic consequence of two ingredients:
(i) a positive feedback between central density and accretion rate
(Bondi), and (ii) a negative feedback between accretion energy and
central density (outflows).  Any sub-grid model with these two
ingredients will produce oscillatory behaviour when coupled to
gravothermal evolution.

\subsection{Observational Prospects}

The SIDM--AGN limit cycle makes several testable predictions:
\begin{enumerate}
  \item An increased scatter in the $\MBH$--$\vdisp$ relation at the
    low-mass end ($\vdisp < 100\kms$), where SIDM cores suppress BH
    growth.
  \item A population of ``over-massive'' BHs in haloes that have
    recently undergone core collapse, detectable as AGN in
    intermediate-mass haloes at $z < 1$.
  \item A redshift-dependent pivot in the $\MBH$--$\vdisp$ relation:
    under-massive at $z > 2$, normal or over-massive at $z < 0.5$.
  \item Quasi-periodic AGN activity with a period of $\sim 0.5$--$1\Gyr$,
    potentially observable in X-ray light curves or relic radio lobes.
\end{enumerate}


% ===================================================================
% 9. CONCLUSIONS
% ===================================================================
\section{Conclusions}
\label{sec:conclusions}

We have presented the first cosmological hydrodynamic simulations that
self-consistently couple SIDM thermodynamics with SMBH growth and AGN
feedback.  Our main conclusions are:

\begin{enumerate}
  \item SIDM core formation ($z \gtrsim 2$) suppresses early SMBH
    growth by reducing the central gas density, with $\MBH$ at $z=2$
    suppressed by factors of $2$--$5$ for $\sigmam =
    5$--$10\,\mathrm{cm^2\,g^{-1}}$.

  \item Gravothermal core collapse ($z \lesssim 1$ for $\sigmam \ge
    5\,\mathrm{cm^2\,g^{-1}}$) triggers a burst of Bondi accretion
    that can cause the SMBH to overshoot the CDM mass by up to a
    factor of $3$.

  \item The coupled DM--SMBH system exhibits a limit cycle:
    gravothermal contraction drives accretion, AGN feedback expels
    baryons and softens the potential, and gas refilling restarts the
    cycle with a period of $\sim 0.5$--$1\Gyr$.

  \item SIDM flattens the $\MBH$--$\vdisp$ relation at the low-mass
    end and increases its scatter, consistent with observational hints
    from dwarf AGN surveys.

  \item Non-elastic SIDM variants differentiate strongly: dissipative
    scattering (\code{dSIDM}, $f_{\rm diss}=0.5$) accelerates core
    collapse by $\sim 2\Gyr$ at fixed nominal $\sigmam$, exothermic
    inelastic scattering (\code{iSIDM}) injects $\delta =
    100\,\mathrm{keV}$ per conversion event and suppresses
    gravothermal collapse entirely, and DM--baryon drag
    (\code{IDM}) at $\sigma_{\rm db}/m =
    0.01\,\mathrm{cm^2\,g^{-1}}$ induces a modest extra baryonic
    infall without measurably altering the dark-matter profile.

  \item AGN feedback extends the gravothermal collapse timescale by
    $30$--$50$ per cent, demonstrating that baryonic physics cannot be
    neglected in predictions of SIDM halo profiles.
\end{enumerate}

These results establish that SMBH growth and SIDM thermodynamics are
inseparable: the dark matter microphysics sets the stage for black
hole accretion, while the black hole's feedback regulates the dark
matter profile.  Future work will extend this study to a statistical
sample of haloes spanning $10^{10}$--$10^{14}\,\Msun$, incorporate
velocity-dependent (Yukawa) cross-sections, and compare directly with
JWST observations of high-redshift AGN.


% ===================================================================
% ACKNOWLEDGEMENTS
% ===================================================================
\section*{Acknowledgements}

This work was supported by the Korea Institute for Advanced Study
(KIAS) grant funded by the Korea government (MSIT).  Numerical
simulations were performed on the KIAS computing cluster.

\section*{Data Availability}

The simulation data underlying this article will be shared on
reasonable request to the corresponding author.  The \curamses\ code
is available at
\url{https://github.com/kjhan0606/cuRAMSES-kjhan}.

% ===================================================================
% REFERENCES
% ===================================================================
\bibliographystyle{mnras}

\begin{thebibliography}{99}

\bibitem[\protect\citeauthoryear{Adhikari et al.}{Adhikari et al.}{2022}]{Adhikari2022}
Adhikari S., et al., 2022, preprint,
\href{https://arxiv.org/abs/2207.10638}{arXiv:2207.10638}

\bibitem[\protect\citeauthoryear{Aubert, Pichon \& Colombi}{Aubert et al.}{2004}]{Aubert2004}
Aubert D., Pichon C., Colombi S., 2004, \href{https://doi.org/10.1111/j.1365-2966.2004.07883.x}{MNRAS, 352, 376}

\bibitem[\protect\citeauthoryear{Balberg, Shapiro \& Inagaki}{Balberg et al.}{2002}]{Balberg2002}
Balberg S., Shapiro S.~L., Inagaki S., 2002, \href{https://doi.org/10.1086/339038}{ApJ, 568, 475}

\bibitem[\protect\citeauthoryear{Bleuler \& Teyssier}{Bleuler \& Teyssier}{2014}]{Bleuler2014}
Bleuler A., Teyssier R., 2014, \href{https://doi.org/10.1093/mnras/stu2005}{MNRAS, 445, 4015}

\bibitem[\protect\citeauthoryear{Boylan-Kolchin, Bullock \& Kaplinghat}{Boylan-Kolchin et al.}{2011}]{Boylan-Kolchin2011}
Boylan-Kolchin M., Bullock J.~S., Kaplinghat M., 2011, \href{https://doi.org/10.1111/j.1745-3933.2011.01074.x}{MNRAS, 415, L40}

\bibitem[\protect\citeauthoryear{Chabrier}{Chabrier}{2003}]{Chabrier2003}
Chabrier G., 2003, \href{https://doi.org/10.1086/376392}{PASP, 115, 763}

\bibitem[\protect\citeauthoryear{Dav{\'e} et al.}{Dav{\'e} et al.}{2019}]{Dave2019}
Dav{\'e} R., Angl{\'e}s-Alc{\'a}zar D., Narayanan D., Li Q., Rafieferantsoa M.~H., Appleby S., 2019, \href{https://doi.org/10.1093/mnras/stz937}{MNRAS, 486, 2827}

\bibitem[\protect\citeauthoryear{de Blok}{de Blok}{2010}]{deBlok2010}
de Blok W.~J.~G., 2010, \href{https://doi.org/10.1155/2010/789293}{Adv.\ Astron., 2010, 789293}

\bibitem[\protect\citeauthoryear{Despali et al.}{Despali et al.}{2022}]{Despali2022}
Despali G., Sparre M., Vegetti S., Vogelsberger M., Zavala J., Marinacci F., 2022, \href{https://doi.org/10.1093/mnras/stab3537}{MNRAS, 510, 2031}

\bibitem[\protect\citeauthoryear{Dubois et al.}{Dubois et al.}{2012}]{Dubois2012}
Dubois Y., Devriendt J., Slyz A., Teyssier R., 2012, \href{https://doi.org/10.1111/j.1365-2966.2011.20236.x}{MNRAS, 420, 2662}

\bibitem[\protect\citeauthoryear{Dubois et al.}{Dubois et al.}{2014}]{Dubois2014}
Dubois Y., et al., 2014, \href{https://doi.org/10.1093/mnras/stu1227}{MNRAS, 444, 1453}

\bibitem[\protect\citeauthoryear{Elbert et al.}{Elbert et al.}{2018}]{Elbert2018}
Elbert O.~D., Bullock J.~S., Garrison-Kimmel S., Rocha M., O\~norbe J., Peter A.~H.~G., 2018, \href{https://doi.org/10.3847/1538-4357/aa9710}{ApJ, 853, 109}

\bibitem[\protect\citeauthoryear{Essig et al.}{Essig et al.}{2019}]{Essig2019}
Essig R., McDermott S.~D., Yu H.-B., Zhong Y.-M., 2019, \href{https://doi.org/10.1103/PhysRevLett.123.121102}{Phys.\ Rev.\ Lett., 123, 121102}

\bibitem[\protect\citeauthoryear{Fry et al.}{Fry et al.}{2015}]{Fry2015}
Fry A.~B., Governato F., Pontzen A., Quinn T., Tremmel M., Anderson L., Menon H., Brooks A.~M., Wadsley J., 2015, \href{https://doi.org/10.1093/mnras/stv1330}{MNRAS, 452, 1468}

\bibitem[\protect\citeauthoryear{Greene et al.}{Greene et al.}{2020}]{Greene2020}
Greene J.~E., Strader J., Ho L.~C., 2020, \href{https://doi.org/10.1146/annurev-astro-032620-021835}{ARA\&A, 58, 257}

\bibitem[\protect\citeauthoryear{Haardt \& Madau}{Haardt \& Madau}{2012}]{Haardt2012}
Haardt F., Madau P., 2012, \href{https://doi.org/10.1088/0004-637X/746/2/125}{ApJ, 746, 125}

\bibitem[\protect\citeauthoryear{Hahn \& Abel}{Hahn \& Abel}{2011}]{Hahn2011}
Hahn O., Abel T., 2011, \href{https://doi.org/10.1111/j.1365-2966.2011.18820.x}{MNRAS, 415, 2101}

\bibitem[\protect\citeauthoryear{Harikane et al.}{Harikane et al.}{2023}]{Harikane2023}
Harikane Y., et al., 2023, \href{https://doi.org/10.3847/1538-4357/ad029e}{ApJ, 959, 39}

\bibitem[\protect\citeauthoryear{Kaplinghat, Tulin \& Yu}{Kaplinghat et al.}{2016}]{Kaplinghat2016}
Kaplinghat M., Tulin S., Yu H.-B., 2016, \href{https://doi.org/10.1103/PhysRevLett.116.041302}{Phys.\ Rev.\ Lett., 116, 041302}

\bibitem[\protect\citeauthoryear{Kim et al.}{Kim et al.}{2024}]{Kim2024}
Kim J., et al., 2024, in preparation

\bibitem[\protect\citeauthoryear{Kim et al.}{Kim et al.}{2026}]{KimInPrep2026}
Kim J., et al., 2026, in preparation

\bibitem[\protect\citeauthoryear{Lee et al.}{Lee et al.}{2021}]{Lee2021}
Lee J., et al., 2021, \href{https://doi.org/10.3847/1538-4357/abd08b}{ApJ, 908, 11}

\bibitem[\protect\citeauthoryear{Snaith et al.}{Snaith et al.}{2018}]{Snaith2018}
Snaith O.~N., Park C., Kim J., Rasera Y., 2018, \href{https://doi.org/10.1093/mnras/sty1731}{MNRAS, 479, 2853}

\bibitem[\protect\citeauthoryear{Tweed et al.}{Tweed et al.}{2009}]{Tweed2009}
Tweed D., Devriendt J., Blaizot J., Colombi S., Slyz A., 2009, \href{https://doi.org/10.1051/0004-6361/200810579}{A\&A, 506, 647}

\bibitem[\protect\citeauthoryear{Vogelsberger et al.}{Vogelsberger et al.}{2019}]{Vogelsberger2019}
Vogelsberger M., Zavala J., Schutz K., Slatyer T.~R., 2019, \href{https://doi.org/10.1093/mnras/sty2970}{MNRAS, 484, 5437}

\bibitem[\protect\citeauthoryear{Kormendy \& Ho}{Kormendy \& Ho}{2013}]{Kormendy2013}
Kormendy J., Ho L.~C., 2013, \href{https://doi.org/10.1146/annurev-astro-082708-101811}{ARA\&A, 51, 511}

\bibitem[\protect\citeauthoryear{Lynden-Bell \& Wood}{Lynden-Bell \& Wood}{1968}]{LyndenBell1968}
Lynden-Bell D., Wood R., 1968, \href{https://doi.org/10.1093/mnras/138.4.495}{MNRAS, 138, 495}

\bibitem[\protect\citeauthoryear{Maiolino et al.}{Maiolino et al.}{2024}]{Maiolino2024}
Maiolino R., et al., 2024, \href{https://doi.org/10.1038/s41586-024-07052-5}{Nature, 627, 59}

\bibitem[\protect\citeauthoryear{Navarro, Frenk \& White}{Navarro et al.}{1996}]{NFW1996}
Navarro J.~F., Frenk C.~S., White S.~D.~M., 1996, \href{https://doi.org/10.1086/177173}{ApJ, 462, 563}

\bibitem[\protect\citeauthoryear{Nishikawa et al.}{Nishikawa et al.}{2020}]{Nishikawa2020}
Nishikawa H., Boddy K.~K., Kaplinghat M., 2020, \href{https://doi.org/10.1103/PhysRevD.101.063009}{Phys.\ Rev.\ D, 101, 063009}

\bibitem[\protect\citeauthoryear{Oh et al.}{Oh et al.}{2015}]{Oh2015}
Oh S.-H., et al., 2015, \href{https://doi.org/10.1088/0004-6256/149/6/180}{AJ, 149, 180}

\bibitem[\protect\citeauthoryear{Oman et al.}{Oman et al.}{2015}]{Oman2015}
Oman K.~A., et al., 2015, \href{https://doi.org/10.1093/mnras/stv1504}{MNRAS, 452, 3650}

\bibitem[\protect\citeauthoryear{Pillepich et al.}{Pillepich et al.}{2018}]{Pillepich2018}
Pillepich A., et al., 2018, \href{https://doi.org/10.1093/mnras/stx2656}{MNRAS, 473, 4077}

\bibitem[\protect\citeauthoryear{Planck Collaboration}{Planck Collaboration}{2016}]{Planck2016}
Planck Collaboration, 2016, \href{https://doi.org/10.1051/0004-6361/201525830}{A\&A, 594, A13}

\bibitem[\protect\citeauthoryear{Planck Collaboration}{Planck Collaboration}{2020}]{Planck2020}
Planck Collaboration, 2020, \href{https://doi.org/10.1051/0004-6361/201833910}{A\&A, 641, A6}

\bibitem[\protect\citeauthoryear{Rasera \& Teyssier}{Rasera \& Teyssier}{2006}]{Rasera2006}
Rasera Y., Teyssier R., 2006, \href{https://doi.org/10.1051/0004-6361:20053116}{A\&A, 445, 1}

\bibitem[\protect\citeauthoryear{Robertson et al.}{Robertson et al.}{2017a}]{Robertson2017a}
Robertson A., Massey R., Eke V., 2017, \href{https://doi.org/10.1093/mnras/stw2670}{MNRAS, 465, 569}

\bibitem[\protect\citeauthoryear{Robertson et al.}{Robertson et al.}{2018}]{Robertson2018}
Robertson A., Massey R., Eke V., Schaye J., Theuns T., 2018, \href{https://doi.org/10.1093/mnrasl/sly024}{MNRAS, 476, L20}

\bibitem[\protect\citeauthoryear{Rocha et al.}{Rocha et al.}{2013}]{Rocha2013}
Rocha M., Peter A.~H.~G., Bullock J.~S., Kaplinghat M., Garrison-Kimmel S., O\~norbe J., Moustakas L.~A., 2013, \href{https://doi.org/10.1093/mnras/sts514}{MNRAS, 430, 81}

\bibitem[\protect\citeauthoryear{Sagunski et al.}{Sagunski et al.}{2021}]{Sagunski2021}
Sagunski L., Gad-Nasr S., Colquhoun B., Robertson A., Tulin S., 2021, \href{https://doi.org/10.1088/1475-7516/2021/01/024}{J.\ Cosmol.\ Astropart.\ Phys., 2021, 024}

\bibitem[\protect\citeauthoryear{Sameie et al.}{Sameie et al.}{2021}]{Sameie2021}
Sameie O., Yu H.-B., Sales L.~V., Vogelsberger M., Zavala J., 2021, \href{https://doi.org/10.1103/PhysRevLett.127.271301}{Phys.\ Rev.\ Lett., 127, 271301}

\bibitem[\protect\citeauthoryear{Santos-Santos et al.}{Santos-Santos et al.}{2020}]{Santos-Santos2020}
Santos-Santos I.~M.~E., et al., 2020, \href{https://doi.org/10.1093/mnras/staa1072}{MNRAS, 495, 58}

\bibitem[\protect\citeauthoryear{Schaye et al.}{Schaye et al.}{2015}]{Schaye2015}
Schaye J., et al., 2015, \href{https://doi.org/10.1093/mnras/stu2058}{MNRAS, 446, 521}

\bibitem[\protect\citeauthoryear{Schutz \& Slatyer}{Schutz \& Slatyer}{2015}]{Schutz2015}
Schutz K., Slatyer T.~R., 2015, \href{https://doi.org/10.1088/1475-7516/2015/01/021}{J.\ Cosmol.\ Astropart.\ Phys., 2015, 021}

\bibitem[\protect\citeauthoryear{Shen et al.}{Shen et al.}{2021}]{Shen2021}
Shen X., Hopkins P.~F., Faucher-Gigu{\`e}re C.-A., Alexander D.~M., Richards G.~T., Ross N.~P., Hickox R.~C., 2021, \href{https://doi.org/10.1093/mnras/stab1992}{MNRAS, 506, 4421}

\bibitem[\protect\citeauthoryear{Smith et al.}{Smith et al.}{2017}]{Smith2017}
Smith B.~D., et al., 2017, \href{https://doi.org/10.1093/mnras/stw3291}{MNRAS, 466, 2217}

\bibitem[\protect\citeauthoryear{Spergel \& Steinhardt}{Spergel \& Steinhardt}{2000}]{Spergel2000}
Spergel D.~N., Steinhardt P.~J., 2000, \href{https://doi.org/10.1103/PhysRevLett.84.3760}{Phys.\ Rev.\ Lett., 84, 3760}

\bibitem[\protect\citeauthoryear{Teyssier}{Teyssier}{2002}]{Teyssier2002}
Teyssier R., 2002, \href{https://doi.org/10.1051/0004-6361:20011817}{A\&A, 385, 337}

\bibitem[\protect\citeauthoryear{Tulin, Yu \& Zurek}{Tulin et al.}{2013}]{Tulin2013}
Tulin S., Yu H.-B., Zurek K.~M., 2013, \href{https://doi.org/10.1103/PhysRevD.87.115007}{Phys.\ Rev.\ D, 87, 115007}

\bibitem[\protect\citeauthoryear{Tulin \& Yu}{Tulin \& Yu}{2018}]{Tulin2018}
Tulin S., Yu H.-B., 2018, \href{https://doi.org/10.1016/j.physrep.2017.11.004}{Phys.\ Rep., 730, 1}

\bibitem[\protect\citeauthoryear{Vogelsberger, Zavala \& Loeb}{Vogelsberger et al.}{2012}]{Vogelsberger2012}
Vogelsberger M., Zavala J., Loeb A., 2012, \href{https://doi.org/10.1111/j.1365-2966.2012.21182.x}{MNRAS, 423, 3740}

\end{thebibliography}

% ===================================================================

\label{lastpage}
\end{document}
